\documentclass[11pt]{article}

\usepackage[T1]{fontenc}

\usepackage[utf8]{inputenc}

\usepackage[margin=1in]{geometry}

\usepackage{graphicx}

\usepackage{amsmath}

\usepackage{amssymb}

\usepackage{booktabs}

\usepackage{array}

\usepackage{tabularx}

\usepackage{caption}

\usepackage{float}

\usepackage{microtype}

\usepackage{mathptmx}

\usepackage{natbib}

\usepackage[hidelinks]{hyperref}

\captionsetup{font=small,labelfont=bf}

\setlength{\parindent}{1.25em}

\setlength{\parskip}{0pt}

\title{Feasibility assessment for estimating refractory epilepsy burden in pediatric Asian populations}

\author{}

\date{}

\begin{document}

\maketitle

\begin{center}\small\textit{Research Memo}\end{center}

\begin{abstract}

This memo assessed whether the prevalence of refractory epilepsy among pediatric populations in Asia could be estimated from the datasets and retrievable sources available in this run. The requested question was evaluated against user-provided dataset hints and a targeted web search for open data related to pediatric epilepsy, refractory epilepsy, Asia, and prevalence. The available materials did not include a usable primary dataset with both pediatric Asian population denominators and refractory-epilepsy status. The retrieved sources were literature and review-style pages rather than a dataset suitable for original quantitative estimation. As a result, no quantitative prevalence estimate was produced, and the run instead documents why the analysis could not be executed faithfully with the available materials.

\end{abstract}

\section{Introduction}
The research question for this run was whether the prevalence of refractory epilepsy among pediatric populations in Asia could be estimated from the datasets and retrievable sources available at hand. This is an important question because it requires a dataset that can support a direct and honest prevalence calculation, ideally with clear case definitions, pediatric age information, geographic relevance to Asia, and a denominator appropriate for estimating refractory disease burden. The validation record indicates that no such artifact-backed result survived review, so the manuscript is structured as a feasibility memo rather than an empirical paper. In that context, the goal is to document the search and assessment process transparently and to explain why no original quantitative result is reported.

\section{Methods}
I inspected the user-provided dataset hints and evaluated whether any were suitable for an original empirical analysis of refractory epilepsy in pediatric Asian populations. BRFSS 2015 Diabetes Health Indicators was rejected because it was unrelated to epilepsy and pediatrics. The NCI Genomic Data Commons API was rejected because it is oncology-focused. MAST Archive and Catalogs was rejected because it is astronomy data. DANDI Archive API was considered potentially relevant to neuroscience, but no specific retrievable dandiset or metadata structure supporting an estimate of refractory-epilepsy prevalence in Asian pediatric populations was identified in this run. I then searched for open data using queries centered on pediatric epilepsy, refractory epilepsy, Asia, and prevalence. The returned materials were literature and review-style resources rather than a directly usable primary dataset. I recorded the search process and the feasibility assessment in the saved artifacts referenced in the ledger. \ref{tab:artifact-1} \ref{tab:artifact-2}

\section{Results}
The search log shows that the queried sources did not yield a usable primary dataset for the requested analysis. The recorded queries included pediatric epilepsy refractory prevalence in Asia, epilepsy children Asia refractory open dataset prevalence, DANDI pediatric epilepsy Asia metadata, and a GDC API query. The log notes that no dataset result was identified, no open patient-level or country-tabulated refractory prevalence dataset was found, and the oncology-focused GDC source was not suitable. \ref{tab:artifact-1} The feasibility notes reach the same conclusion: BRFSS was unrelated, GDC was unsuitable, MAST was irrelevant, and DANDI did not provide a specific retrievable dataset supporting the requested estimate. The returned web materials were literature pages and reviews, not an accessible primary dataset for original estimation. \ref{tab:artifact-2} No quantitative prevalence result was produced.

\begin{table}[tbp]
\centering
\caption{Structured log of the search queries used in this run and why they did not yield a usable primary dataset for the requested analysis.}
\label{tab:artifact-1}
\begin{tabular}{>{\raggedright\arraybackslash}p{0.15\linewidth}>{\raggedright\arraybackslash}p{0.15\linewidth}}
\toprule
query & note \\
\midrule
pediatric epilepsy refractory prevalence Asia study open data dataset & No dataset result identified in provided searchable results; results were litera \\
epilepsy children Asia refractory open dataset prevalence & No open patient-level or country-tabulated refractory prevalence dataset found i \\
DANDI pediatric epilepsy asia metadata & Relevant archive exists for neuroscience data, but no evidence from returned res \\
site:api.gdc.cancer.gov epilepsy pediatric asia & Not relevant; GDC is oncology-focused and unsuitable for epilepsy prevalence que \\
\bottomrule
\end{tabular}
\end{table}

\begin{table}[tbp]
\centering
\caption{Narrative feasibility assessment explaining why the provided datasets and retrieved sources could not support an honest original quantitative study of refractory epilepsy in pediatric Asian populations.}
\label{tab:artifact-2}
\begin{tabular}{>{\raggedright\arraybackslash}p{0.15\linewidth}}
\toprule
This run evaluated the user question 'Epilepsy in pediatric Asian populations ho \\
\midrule
- BRFSS 2015 Diabetes Health Indicators: unrelated to epilepsy or pediatrics. \\
- NCI Genomic Data Commons API: oncology metadata \\
- MAST Archive and Catalogs: astronomy data \\
- DANDI Archive API: neuroscience archive potentially relevant to epilepsy recor \\
I also searched the web for open data specifically on pediatric epilepsy refract \\
\bottomrule
\end{tabular}
\end{table}

\section{Discussion}
The evidence from this run supports a narrow but important conclusion: the requested prevalence question could not be answered honestly with the available materials. The central obstacle was not a failed computation but the absence of a fit-for-purpose dataset. To estimate refractory epilepsy burden in pediatric Asian populations, one would need a source with explicit pediatric case data, a clear definition of refractory epilepsy, and a denominator that permits prevalence estimation in an Asian population. The search process did not surface such a source, and the materials that were found were not sufficient for original quantitative analysis. The decision to report no numeric result is consistent with the goal of avoiding overinterpretation of secondary literature or unrelated datasets. The memo therefore serves as a documented feasibility assessment and a record of the methods used in this run.

\section{Limitations}
No provided dataset was suitable for estimating refractory epilepsy prevalence in pediatric Asian populations. The web search results available in this run surfaced literature summaries and reviews, not a directly retrievable primary dataset for original quantitative estimation. Without a usable cohort, registry, survey, or case-series dataset containing both pediatric Asian population denominators and refractory-epilepsy status, meaningful original analysis was not possible. Because the requested study could not be executed faithfully with available materials, no quantitative results are reported.

\bibliographystyle{plainnat}
\bibliography{references}

\end{document}
